\chapter*{Annexe}
\addcontentsline{toc}{chapter}{Annexe}

\section*{Annexe 1 : Modèle Relationnel de Données}
Le stockage des informations repose sur un modèle relationnel standard. Les principales tables sont les suivantes :

\begin{itemize}
    \item \textbf{Tickets} : Cette table centrale stocke toutes les réclamations, avec leur sujet, leur description, et les clés étrangères vers le client et le statut.
    \item \textbf{Utilisateurs} : Contient les informations de connexion (email, mot de passe haché) et les profils (Client, RMQ, Admin).
    \item \textbf{Statuts} : Définit le cycle de vie d'un ticket (Nouveau -> En cours -> Clôturé).
    \item \textbf{Historique} : Permet une traçabilité complète de chaque action effectuée sur un ticket (Qui a fait quoi et quand).
\end{itemize}
